% 量子力学1 量子力学1-11 問題8（139-parity） 解いた日 2026-11-20
%% kotae
\toi{1}{◎}{}
\begin{align*}
\hat{P}^{2}\psi(x)&=\hat{P}\bigl(\hat{P}\psi(x)\bigr)\\
&=\hat{P}\psi(-x)\\
&=\psi(x)\qquad\therefore\ \hat{P}^{2}=\hat{1}
\end{align*}
$\hat{P}$ の固有値を $\lambda$，固有関数を $\chi(x)$ とする。
\[ \hat{P}\chi(x)=\lambda\chi(x)\qquad\text{両辺に $\hat{P}$ を作用させる。} \]
\begin{align*}
&\hat{P}^{2}\chi(x)=\lambda\hat{P}\chi(x)\\
&\hat{1}\chi(x)=\lambda^{2}\chi(x)\\
&\therefore\ \lambda^{2}=1\qquad\lambda=\pm1
\end{align*}
$\lambda=1$ のとき、
\[ \hat{P}\chi(x)=\chi(-x)=\chi(x) \]
$\therefore$ $\chi(x)$ は偶パリティ

$\lambda=-1$ のとき、
\[ \hat{P}\chi(x)=\chi(-x)=-\chi(x) \]
$\chi(x)$ は奇パリティ
\toi{2}{◎}{}
\[ \int\psi^{*}\hat{P}\psi\dd x=\int\hat{P}\psi^{*}\psi\dd x\quad\text{を示す} \]
\begin{align*}
\text{左辺}&=\int_{-\infty}^{\infty}\psi^{*}(x)\psi(-x)\dd x\qquad y=-x\ \text{とする。}\\
&=\int_{\infty}^{-\infty}\psi^{*}(-y)\psi(y)\cdot{-\dd y}\\
&=\int_{-\infty}^{\infty}\psi^{*}(-y)\psi(y)\dd y\\
\text{右辺}&=\int_{-\infty}^{\infty}\psi^{*}(-x)\psi(x)\dd x
\end{align*}
より、左辺＝右辺が成り立つ。

よって $\hat{P}$ はエルミート演算子
\toi{3}{◎}{}
\begin{align*}
\hat{P}\hat{x}\psi(x)&=\hat{P}\bigl(x\psi(x)\bigr)\\
&=-x\psi(-x)\\
-\hat{x}\hat{P}\psi(x)&=-x\psi(-x)\quad\text{より、}\ \hat{P}\hat{x}=-\hat{x}\hat{P}
\end{align*}
\begin{align*}
\hat{P}\hat{p}\psi(x)&=\hat{P}\bigl(-\ii\hbar\psi'(x)\bigr)\\
&=-\ii\hbar\psi'(-x)\\
\hat{p}\hat{P}\psi(x)&=-\ii\hbar\dv{}{x}\bigl(\psi(-x)\bigr)\\
&=\ii\hbar\psi'(-x)\qquad\text{よって}\ \hat{P}\hat{p}=-\hat{p}\hat{P}
\end{align*}
\[ \comm{\hat{P}}{\hat{H}}=\frac{1}{2m}\comm{\hat{P}}{\hat{p}^{2}}+\comm{\hat{P}}{V(x)} \]
\begin{align*}
\comm{\hat{P}}{\hat{p}^{2}}&=\comm{\hat{P}}{\hat{p}}\hat{p}+\hat{p}\comm{\hat{P}}{\hat{p}}\\
&=(\hat{P}\hat{p}-\hat{p}\hat{P})\hat{p}+\hat{p}(\hat{P}\hat{p}-\hat{p}\hat{P})\\
&=2\hat{P}\hat{p}^{2}+2\hat{p}\hat{P}\hat{p}\\
&=2\hat{P}\hat{p}^{2}-2\hat{P}\hat{p}^{2}\\
&=0
\end{align*}
\begin{align*}
\comm{\hat{P}}{V(x)}\psi(x)&=\hat{P}V(x)\psi(x)-V(x)\hat{P}\psi(x)\\
&=V(-x)\psi(-x)-V(x)\psi(-x)\\
&=\bigl(V(-x)-V(x)\bigr)\hat{P}\psi(x)
\end{align*}
\[ \comm{\hat{P}}{V(x)}=\bigl(V(-x)-V(x)\bigr)\hat{P} \]
よって $\comm{\hat{P}}{\hat{H}}=0$ となるのは $V(-x)=V(x)$ のときのみ
\toi{4}{○}{記述}
\begin{align*}
\expval{\hat{P}}&=\int_{-\infty}^{\infty}\psi^{*}(x)\hat{P}\psi(x)\dd x\\
&=\int_{-\infty}^{\infty}\psi(x)\psi(-x)\dd x\\
&=(\pi\sigma^{2})^{-\frac{1}{2}}\int_{-\infty}^{\infty}\exp\biggl(\frac{-(x-x_{0})^{2}}{2\sigma^{2}}\\
&\qquad+\frac{-(-x-x_{0})^{2}}{2\sigma^{2}}\biggr)\dd x
\end{align*}
（矢印で指数の部分につなぐ）
\begin{align*}
&\frac{1}{2\sigma^{2}}\Bigl(-(x^{2}-2x_{0}x+x_{0}^{2})-(x^{2}+2x_{0}x+x_{0}^{2})\Bigr)\\
&=\frac{-(x^{2}+x_{0}^{2})}{\sigma^{2}}
\end{align*}
\begin{align*}
&\int_{-\infty}^{\infty}\exp\left(-\frac{x_{0}^{2}}{\sigma^{2}}\right)\exp\left(-\frac{x^{2}}{\sigma^{2}}\right)\dd x\\
&=\exp\left(-\frac{x_{0}^{2}}{\sigma^{2}}\right)\sqrt{\pi\sigma^{2}}
\end{align*}
\[ \expval{\hat{P}}=\exp\left(-\frac{x_{0}^{2}}{\sigma^{2}}\right) \]
\[ \psi(x)=\underbrace{\psi_{+}(x)}_{\text{偶}}+\underbrace{\psi_{-}(x)}_{\text{奇}} \]
\begin{align*}
&\hat{P}\psi(x)=\psi_{+}(x)-\psi_{-}(x)\\
&(1+\hat{P})\psi(x)=2\psi_{+}(x)\\
&\therefore\ \psi_{+}(x)=\frac{1+\hat{P}}{2}\psi(x)\\
&(1-\hat{P})\psi(x)=2\psi_{-}(x)\\
&\therefore\ \psi_{-}(x)=\frac{1-\hat{P}}{2}\psi(x)
\end{align*}
\begin{align*}
\int_{-\infty}^{\infty}\abs{\psi_{+}(x)}^{2}\dd x&=\frac{1}{4}\int\psi^{2}(x)+2\psi(x)\psi(-x)\\
&\qquad+\psi^{2}(-x)\dd x\\
&=\frac{1}{4}\bigl(2+2\expval{\hat{P}}\bigr)\\
&=\frac{1}{2}\left(1+\exp\left(-\frac{x_{0}^{2}}{\sigma^{2}}\right)\right)\\
\int_{-\infty}^{\infty}\abs{\psi_{-}(x)}^{2}\dd x&=\frac{1}{4}\int\psi^{2}(x)-2\psi(x)\psi(-x)\\
&\qquad+\psi^{2}(-x)\dd x\\
&=\frac{1}{4}\bigl(2-2\expval{\hat{P}}\bigr)\\
&=\frac{1}{2}\left(1-\exp\left(-\frac{x_{0}^{2}}{\sigma^{2}}\right)\right)
\end{align*}
\naze{$\int\abs{\psi_{\pm}}^{2}\dd x$ がパリティを測って $\pm1$ を得る確率になる根拠が無い。$\psi_{+}$ と $\psi_{-}$ は $\hat{P}$ の異なる固有値の固有関数で，$\hat{P}$ がエルミート（(2)）なので直交する。だから $\int\abs{\psi}^{2}=\int\abs{\psi_{+}}^{2}+\int\abs{\psi_{-}}^{2}$ と分かれ，各成分の積分が確率になる。}
%% kizuki
\item \mae{第7章 問題4・11-07} 確率と読める理由を書かなかった（前回は $\abs{\tilde{\psi}}^{2}\dd k$ が確率になる理由）。測定で固有値 $\pm1$ を得る確率は，$\psi$ を固有関数どうしの直交する成分に分けたときの各成分の規格化の積分である。直交していなければ $\int\abs{\psi}^{2}=\int\abs{\psi_{+}}^{2}+\int\abs{\psi_{-}}^{2}$ が成り立たず確率と読めない。(2) のエルミート性がここで使われる。
\item 確かめ：$P_{+}+P_{-}=1$，$x_{0}=0$（偶関数）で $P_{+}=1$ になる。
